\labeledsection{Unary Natural Numbers}{sec:unarynatnums}
\definition{Unary Natural Numbers}{
    Terms generated by the following: $u ::= 0 \mid s(u)$ where $0$ denotes zero and $s(u)$ denotes the successor of $u$. i.e. $\set{0, s(0), s(s(0)), s(s(s(0))), \cdots}$.
}{unarynatnums}

\definition{Piano Axioms}{
    The following set of \linkterm{axioms}{axiom} are called the Piano axioms. They can be used to \textit{reason} about natural numbers:
}{piano}
\begin{itemize}
    \item \textbf{P1.} $0$ is a \linkterm{unary natural number}{unarynatnums}.
    \item \textbf{P2.} Every \linkterm{unary natural number}{unarynatnums} has a successor that is \linkterm{unary natural number}{unarynatnums} and that is different from it.
    \item \textbf{P3.} $0$ is not a successor of any \linkterm{unary natural number}{unarynatnums}.
    \item \textbf{P4.} Different \linkterm{unary natural numbers}{unarynatnums} have different successors.
    \item \textbf{P5.} \refterm{Induction Axiom}{induction_axiom}. Every \linkterm{unary natural number}{unarynatnums} possess a property $P$ if
    \begin{itemize}
        \item $0$ has the property $P$, and
        \item the successor of every \linkterm{unary natural number}{unarynatnums} that has the property $P$ also possesses property $P$
    \end{itemize}
\end{itemize}

More formally, we call the \linkterm{set}{def:set} of \linkterm{unary natural numbers}{unarynatnums} $\mathbb{N}_1$, and we say $P(n)$ if a \linkterm{unary natural number}{unarynatnums} $n \in \mathbb{N}_1$ has a property $P$. Therefore, we express the axioms as follows:
\begin{itemize}
\item $0 \in \mathbb{N}_1$
\item $\forall n \in \mathbb{N}_1.s(n) \in \mathbb{N}_1 \land n \neq s(n)$ 
\item $\neg (\exists n \in \mathbb{N}_1.0=s(n))$
\item $\forall n \in \mathbb{N}_1. \forall m \in \mathbb{N}_1. n \neq m \Rightarrow s(n) \neq s(m)$
\item $\forall P.P(0) \land (\forall n \in \mathbb{N}_1.P(n) \Rightarrow P(s(n))) \Rightarrow (\forall m \in \mathbb{N}_1.P(m))$
\end{itemize}


\Theorem{Every \linkterm{unary natural number}{unarynatnums} is either zero ($0$) or the successor of a \linkterm{unary natural number}{unarynatnums}.}{th:zeroorsucc}

\Proof{
    We make use of the induction axiom P5: We use the property $P$ of "being zero or a successor" and prove the statement by convincing ourselves of the prerequisites of:
    \begin{enumerate}
        \item $0$ is zero so $0$ is "zero or a successor".
        \item Let $n$ be arbitrary \linkterm{unary natural number}{unarynatnums} that is "zero or a successor".
        \item Then its successor is a "successor", so the successor of $n$ is also "zero or a successor".
        \item Since we have taken $n$ arbitrary we have shown that for any $n$, its successor has property $P$.
        \item Property $P$ holds for all \linkterm{unary natural number}{unarynatnums} by P5, so we have proven the assertion.
    \end{enumerate}
}

\definition{Mathematical Induction}{
    A method for proving a \linkterm{statement}{math_statement} $P(n)$ is true for every natural number $n$. A \linkterm{proof}{proofs} by mathematical induction consists of two steps:
    \begin{itemize}
        \item The \refterm{base case}{base_case}: Prove that $P(0)$ is true.
        \item The \refterm{induction step}{induction_step}: Prove that for every $n$, if $P(n)$ is true then $P(n+1)$ is also true.
    \end{itemize}
    The hypothesis in the \linkterm{induction step}{induction_step}, that $P(n)$ is true for a particular $n$, is called the \refterm{induction hypothesis}{induction_hypothesis}
}{mathematical_induction}

\definition{The addition "function"}{
    We define the addition operation $+$ procedurally:
    \begin{itemize}
        \item Adding $0$ to a number does not change it, expressed as: $n + 0 = n$.
        \item Adding $m$ to the successor of $n$ yields the successor of $m + n$, expressed as: $m + s(n) = s(m + n)$.
    \end{itemize}
}{addition}

\Theorem{For all \linkterm{unary natural number}{unarynatnums} $n, m, l$, we have $(n + m) + l = n + (m + l)$.}{th:associativeaddition}
\Proof{
    We prove this by induction on $l$.
    
    Let $P(l)$ be the property: 
    \[
    (n + m) + l = n + (m + l).
    \]
    \begin{itemize}
        \item \linkterm{base case}{base_case}: 
        \[
            (n + m) + 0 = (n + m) = n + (m + 0)
        \]
        \item \linkterm{step case}{induction_step}: 
        Assume $(n + m) + l = n + (m + l)$ for an arbitrary $l$. We must show 
        \[
            (n + m) + s(l) = n + (m + s(l)).
        \]
        \begin{align*}
            (n + m) + s(l) 
                &= s\big((n + m) + l\big) \\
                &\overset{\text{IH}}{=} s\big(n + (m + l)\big) \\
                &= n + s(m + l) \\
                &= n + (m + s(l)).
        \end{align*}
        Hence the property holds for $s(l)$.
    \end{itemize}
}

\definition{Multiplication}{
    The \textbf{unary multiplication} operation can be defined by the equations:
    \begin{itemize}
        \item $n \cdot 0 = 0$, and
        \item $n \cdot s(m) = n + (n \cdot m)$
    \end{itemize}

}{multiplication}

\definition{Exponentiation}{
    The \textbf{unary exponentiation} operation can be defined by the equations:
    \begin{itemize}
        \item $exp(n, 0) = s(0)$, and
        \item $exp(n,s(m)) = n \cdot exp(n,m)$
    \end{itemize}
}{exponentiation}

\definition{Summation}{
    The \textbf{unary summation} operation can be defined by the equations:
    \begin{itemize}
        \item $\bigoplus_{i=0}^0 (n_i) = 0$, and
        \item $\bigoplus_{i=0}^{s(m)}(n_i)=n_{s(m)} \bigoplus \bigoplus_{i=0}^m (n_i)$
    \end{itemize}
}{summation}

Example:
\begin{align*}
\bigoplus_{i=0}^{s(s(0))}(n_i)&=n_{s(s(0))} \bigoplus \bigoplus_{i=0}^{s(0)} (n_i) \\
&=n_{s(s(0))} \bigoplus n_{s(0)} \bigoplus \bigoplus_{i=0}^0 (n_i)\\
&=n_{s(s(0))} \bigoplus n_{s(0)} \bigoplus 0\\
\end{align*}

\definition{Product}{
    The \textbf{unary product} operation can be defined by the equations:
    \begin{itemize}
        \item $\bigodot_{i=0}^0 (n_i) = s(0)$, and
        \item $\bigodot_{i=0}^{s(m)}(n_i)=n_{s(m)} \bigodot \bigodot_{i=0}^m (n_i)$
    \end{itemize}
}{product}